% Options for packages loaded elsewhere
\PassOptionsToPackage{unicode}{hyperref}
\PassOptionsToPackage{hyphens}{url}
\PassOptionsToPackage{dvipsnames,svgnames,x11names}{xcolor}
\documentclass[
  11pt]{article}
\usepackage{xcolor}
\usepackage[margin=1in]{geometry}
\usepackage{amsmath,amssymb}
\setcounter{secnumdepth}{-\maxdimen} % remove section numbering
\usepackage{iftex}
\ifPDFTeX
  \usepackage[T1]{fontenc}
  \usepackage[utf8]{inputenc}
  \usepackage{textcomp} % provide euro and other symbols
\else % if luatex or xetex
  \usepackage{unicode-math} % this also loads fontspec
  \defaultfontfeatures{Scale=MatchLowercase}
  \defaultfontfeatures[\rmfamily]{Ligatures=TeX,Scale=1}
\fi
\usepackage{lmodern}
\ifPDFTeX\else
  % xetex/luatex font selection
\fi
% Use upquote if available, for straight quotes in verbatim environments
\IfFileExists{upquote.sty}{\usepackage{upquote}}{}
\IfFileExists{microtype.sty}{% use microtype if available
  \usepackage[]{microtype}
  \UseMicrotypeSet[protrusion]{basicmath} % disable protrusion for tt fonts
}{}
\makeatletter
\@ifundefined{KOMAClassName}{% if non-KOMA class
  \IfFileExists{parskip.sty}{%
    \usepackage{parskip}
  }{% else
    \setlength{\parindent}{0pt}
    \setlength{\parskip}{6pt plus 2pt minus 1pt}}
}{% if KOMA class
  \KOMAoptions{parskip=half}}
\makeatother
\ifLuaTeX
\usepackage[bidi=basic]{babel}
\else
\usepackage[bidi=default]{babel}
\fi
\babelprovide[main,import]{american}
% get rid of language-specific shorthands (see #6817):
\let\LanguageShortHands\languageshorthands
\def\languageshorthands#1{}
\ifLuaTeX
  \usepackage[english]{selnolig} % disable illegal ligatures
\fi
\setlength{\emergencystretch}{3em} % prevent overfull lines
\providecommand{\tightlist}{%
  \setlength{\itemsep}{0pt}\setlength{\parskip}{0pt}}
\AtBeginDocument{\hypersetup{pdfhighlight=/N}}
\AtBeginDocument{\special{pdf:trailerid [<f44dce828c4fc99aea778d20e567782a> <f44dce828c4fc99aea778d20e567782a>]}}
\usepackage{amsmath,amssymb,amsthm,mathtools}
\usepackage{microtype}
\usepackage{titling}
\setlength{\droptitle}{-1.5em}
\pretitle{\begin{center}\LARGE}
\posttitle{\par\end{center}\vspace{0.25em}}
\preauthor{}
\postauthor{}
\predate{\begin{center}\normalsize}
\postdate{\par\end{center}\vspace{-0.25em}}
\usepackage{needspace}
\let\papersection\subsection
\renewcommand{\subsection}{\Needspace{10\baselineskip}\papersection}
\let\papersubsection\subsubsection
\renewcommand{\subsubsection}{\Needspace{7\baselineskip}\papersubsection}
\raggedbottom
\frenchspacing
\hyphenation{entirely}
\clubpenalty=10000
\widowpenalty=10000
\displaywidowpenalty=10000
\usepackage{enumitem}
\setlist[enumerate]{
  before=\Needspace{12\baselineskip},
  midpenalty=10000,
  endpenalty=5000
}
\usepackage{fancyhdr}
\pagestyle{fancy}
\fancyhf{}
\fancyhead[L]{Dean's conjecture at five}
\fancyfoot[C]{\thepage}
\setlength{\headheight}{14pt}
\usepackage{bookmark}
\IfFileExists{xurl.sty}{\usepackage{xurl}}{} % add URL line breaks if available
\urlstyle{same}
\hypersetup{
  pdftitle={Dean's conjecture at five},
  pdfauthor={Kamil Sieniawski},
  pdflang={en-US},
  colorlinks=true,
  linkcolor={black},
  filecolor={Maroon},
  citecolor={black},
  urlcolor={black},
  pdfcreator={LaTeX via pandoc}}

\title{Dean's conjecture at five}
\author{}
\date{28 July 2026}

\begin{document}
\maketitle

\subsection{Abstract}\label{abstract}

We prove that every nonempty finite simple graph of minimum degree at
least five contains a cycle whose length is divisible by five. The proof
is by contradiction. The two-connected case is ruled out by a theorem of
Bai, Grzesik, Li, and Prorok, so we pass to a nontrivial end block \(B\)
with cut vertex \(c\) and set \(J=B-c\). At most four vertices of \(J\)
have lost one unit of degree, and all are adjacent to \(c\). After
proving that \(J\) is 2-connected and then 3-connected, we repeatedly
reintroduce \(c\) in rooted auxiliary graphs to repair precisely these
degree deficits. This \emph{root lifting} forces 4-connectivity and
excludes \(K_4^-\), triangles, and four-cycles. Thus \(J\) has girth at
least six. A minimum theta, together with its exterior, then provides
two three-path families whose residues yield the final contradiction
modulo five. Together with the previously known cases, this establishes
Dean's conjecture for every \(k\ge3\).

\textbf{Formal verification.} The modulus-five theorem is verified in
Lean. The paper and Lean sources are available at
\href{https://github.com/pkchv/dean-conjecture}{github.com/pkchv/dean-conjecture}.
Section 1.3 explains the relationship between the proof presented here,
the published results it uses, and the formal proof.

\textbf{AI disclosure.} The mathematical proof, manuscript, and Lean
formalization were generated entirely by GPT-5.6 Sol and GPT-5.6 Sol
Pro.

\newpage

\subsection{1. Introduction, proof architecture, and published
results}\label{introduction-proof-architecture-and-published-results}

Dean conjectured that, for every integer \(k\ge3\), every nonempty
finite simple graph of minimum degree at least \(k\) contains a cycle
whose length is divisible by
\(k\).\textsuperscript{\hyperref[ref-dean-1991]{{[}5{]}}} The cases
\(k=3\) and \(k=4\) were proved by Chen and
Saito\textsuperscript{\hyperref[ref-chen-saito-1994]{{[}2{]}}} and Dean,
Lesniak, and
Saito\textsuperscript{\hyperref[ref-dean-lesniak-saito-1993]{{[}6{]}}},
respectively. Luo, Ma, and Zhao recently resolved every case
\(k\ge6\).\textsuperscript{\hyperref[ref-luo-ma-zhao-2026]{{[}8{]}}}
Thus \(k=5\) was the unique remaining case, as also recorded by Choi,
Chu, Kim, and
Park.\textsuperscript{\hyperref[ref-choi-chu-kim-park-2026]{{[}4{]}}}
The theorem proved here therefore completes Dean's conjecture.

For brevity, write BGLP for Bai--Grzesik--Li--Prorok, COY for
Chiba--Ota--Yamashita, and GHLM for Gao--Huo--Liu--Ma.

\subsubsection{1.1. Proof architecture}\label{proof-architecture}

The proof is by contradiction. BGLP Theorem 1.3, together with direct
cycles in its two exceptional families, disposes of a two-connected
counterexample. We therefore work in a nontrivial end block \(B\) with
cut vertex \(c\). Every vertex of \(B-c\) has degree at least five in
\(B\), while \(2\le d_B(c)\le4\). Hence, for \(J=B-c\), every vertex has
degree at least four, and the degree-four vertices form a set \[
 D\subseteq N_B(c),\qquad |D|\le4.
\] In the auxiliary graphs constructed below, after the prescribed
boundary losses are accounted for, precisely these vertices can remain
one degree short of a rooted admissible-path theorem.

Under these standing hypotheses, the structural reduction follows the
chain \[
 J\text{ 2-connected}
 \ \Longrightarrow\ \ J\text{ 3-connected}
 \ \Longrightarrow\ \ J\text{ 4-connected}
 \ \Longrightarrow\ \operatorname{girth}(J)\ge6.
\] The first implication uses BGLP's connectivity lemma. For the second
and for the local exclusions that follow, the recurring device is root
lifting: when there is no deficient vertex we apply GHLM directly, and
when there is one we use it as the COY exception. When there are at
least two, we adjoin \(c\) and restore one edge at each; the COY theorem
then permits \(c\) as its single exceptional vertex. This yields
admissible path families on opposite sides of a separator or outside a
small core; joining those families produces five admissible cycles and
hence a cycle of length \(0\) modulo five.

Once triangles and four-cycles are excluded, the standing counterexample
assumption also excludes 5- and 10-cycles, so \(J\) has girth at least
six and no 10-cycle. Choose a minimum-order theta \(H\). GHLM Lemma 5.10
places \(H\) in a small induced envelope \(T\); the elementary
subdivision-attachment lemma proved in Section 7 then shows that every
vertex outside \(T\) has at most one neighbor in \(T\). We prove that
the exterior \(M=J-V(T)\) is 2-connected. Lemma 2.4 supplies three
\(x,y\)-paths in \(T\) with distinct residues modulo five, while root
lifting supplies three admissible paths through \(M\), possibly using
\(c\). The resulting nine cycles contradict the \(3+3\) residue lemma,
Lemma 2.3.

\subsubsection{1.2. Notation}\label{notation}

All graphs are finite, undirected, and simple; any graph subject to a
minimum-degree hypothesis is nonempty. Paths are simple unless
explicitly called walks, and length means number of edges. For a graph
\(G\), write \(d_G(v)\) and \(N_G(v)\) for the degree and neighborhood
of \(v\), and write \(\delta(G)\) and \(\kappa(G)\) for the minimum
degree and vertex connectivity. If \(X\subseteq V(G)\), put \[
 d_X(v)=|N_G(v)\cap X|,
 \qquad
 N_X(Q)=N_G(Q)\cap X
\] when the ambient graph is clear, where
\(N_G(Q)=\left(\bigcup_{v\in Q}N_G(v)\right)\setminus Q\). Thus
\(d_H(v)\) denotes degree in the graph \(H\), whereas \(d_X(v)\) counts
ambient neighbors in the vertex set \(X\).

The graph \(G-X\) is obtained by deleting \(X\), and \(G[X]\) is induced
by \(X\). The notation \(G-xy\) deletes the edge \(xy\) if present and
otherwise leaves \(G\) unchanged; similarly, \(G+xy\) adds the edge if
absent and otherwise leaves \(G\) unchanged. A \(k\)-vertex cut has
exactly \(k\) distinct vertices. Two edges are \textbf{independent} if
they have no common endpoint. We use the standard block convention that
a 2-connected graph is its unique end block and that all its vertices
are inner vertices.

A \textbf{theta} is the union of three distinct, internally
vertex-disjoint paths with the same two distinct ends. A sequence of
paths or cycles is \textbf{admissible} if its lengths form an arithmetic
progression with common difference one or two; for paths, the shortest
length is at least two.

We use without further comment that a 2-connected graph which is not a
cycle contains a theta, that an induced theta is 2-connected and has
maximum degree at most three, and that adjoining a vertex with two
distinct neighbors to a 2-connected graph preserves 2-connectivity. The
required forms of the block-cut-tree argument, edge contraction, and
Menger's theorem appear where they are used.

\subsubsection{1.3. Published results used in the modulus-five
proof}\label{published-results-used-in-the-modulus-five-proof}

The modulus-five proof uses the following five published results.

\textbf{Rooted admissible-path theorem (GHLM, Theorem
3.1).}\textsuperscript{\hyperref[ref-gao-huo-liu-ma-2022]{{[}7{]}}} Let
\(q\) be a positive integer, and let \(A\) have distinct roots \(x,y\)
such that \(A+xy\) is 2-connected. If \[
 d_A(v)\ge q+1\qquad(v\in V(A)\setminus\{x,y\}),
\] then \(A\) contains \(q\) admissible \(x,y\)-paths.

\textbf{One-exception rooted theorem (the specialization of COY, Theorem
3, stated as BGLP, Theorem
2.2).}\textsuperscript{\hyperref[ref-chiba-ota-yamashita-2023]{{[}3{]}},~\hyperref[ref-bglp-2026]{{[}1{]}}}
Let \(q\) be a positive integer, let \(|V(A)|\ge4\), and let
\(x,y,z\in V(A)\), with \(x\ne y\), such that \(A+xy\) is 2-connected.
If \[
 d_A(v)\ge q+1\qquad(v\in V(A)\setminus\{x,y,z\}),
\] then \(A\) contains \(q\) admissible \(x,y\)-paths. No distinctness
condition is imposed on \(z\) relative to the roots.

\textbf{Connectivity lemma (BGLP, Lemma
2.3).}\textsuperscript{\hyperref[ref-bglp-2026]{{[}1{]}}} Let \(k\ge2\)
and \(r\ge0\) be integers. If a 2-connected graph \(H\) has minimum
degree at least \(k-r\) and has no \(k\) admissible cycles, then \[
 \kappa(H)\ge \frac{k+2}{2}-r.
\]

\textbf{Two-connected minimum-degree theorem (BGLP, Theorem
1.3).}\textsuperscript{\hyperref[ref-bglp-2026]{{[}1{]}}} Let \(H\) have
order \(n\). If \(H\) is 2-connected and \(\delta(H)\ge k\ge4\), then
\(H\) has \(k\) admissible cycles, unless it is isomorphic to one of the
following two graphs: \[
 H\cong K_{k+1}
 \qquad\text{or}\qquad
 H\cong K_{k,n-k}.
\]

\Needspace{8\baselineskip}

\textbf{Minimum-theta structure (GHLM, Lemma
5.10).}\textsuperscript{\hyperref[ref-gao-huo-liu-ma-2022]{{[}7{]}}} Let
\(A\) have girth at least six and no 10-cycle, and let \(H\) be a
minimum-order theta in \(A\). Then \(H\) is induced; at most one vertex
of \(A-V(H)\) has at least two neighbors in \(H\); and, if such a vertex
\(v\) exists, then \(A[V(H)\cup\{v\}]\) is isomorphic to a
one-subdivision of \(K_4\).

The auxiliary forms of the two-link theorem, the subdivision-attachment
lemma, and the complete-bipartite path ranges are proved where they are
used and are also formalized in Lean. Lemma 2.4 gives the remaining
theta-residue calculation.

The proof uses the five published results above as stated. The Lean
formalization proves internally the specializations used here: its
rooted-path interfaces cover every value \(2\le q\le4\) occurring below,
and its two-connected branch applies the Sections 4--7 argument with
empty deficiency rather than formalizing BGLP Theorem 1.3 verbatim.

\subsection{2. Length arithmetic and simple-cycle
assembly}\label{length-arithmetic-and-simple-cycle-assembly}

The following lemmas contain all arithmetic used to combine path
families.

\textbf{Lemma 2.1 (sumsets of admissible progressions).} Let \[
 \mathcal A=\{a+i\alpha:0\le i<s\},\qquad
 \mathcal B=\{b+j\beta:0\le j<t\},
\] where \(a,b\in\mathbb Z\), \(s,t\) are positive integers, and
\(\alpha,\beta\in\{1,2\}\). Then \(\mathcal A+\mathcal B\) contains an
arithmetic progression of common difference one or two with at least
\(s+t-1\) terms.

\textbf{Proof.} If \(\alpha=\beta\), the entire sumset is \[
 \{a+b+h\alpha:0\le h\le s+t-2\}.
\] If \(\{\alpha,\beta\}=\{1,2\}\), translate and, if necessary,
interchange the pairs \((\mathcal A,s)\) and \((\mathcal B,t)\) so that
\[
 \mathcal A=\{0,1,\ldots,s-1\},\qquad
 \mathcal B=\{0,2,\ldots,2(t-1)\}.
\] If \(s=1\), the sum is the \(t\)-term progression \(\mathcal B\). If
\(s\ge2\), the intervals \([2j,2j+s-1]\), \(0\le j<t\), overlap or are
adjacent. Their union is therefore an interval of at least \(s+t-1\)
integers. \(\square\)

\textbf{Lemma 2.2 (concatenation).} Let \(H\) be a subgraph of \(G\),
and let \(X,Y\subseteq V(H)\) be disjoint. Let \[
 \mathcal A=\{a_1,\ldots,a_s\},\qquad
 \mathcal B=\{b_1,\ldots,b_t\}
\] be increasing arithmetic progressions of positive integers, so that
\(a_i=a_1+(i-1)\alpha\) and \(b_j=b_1+(j-1)\beta\) for
\(\alpha,\beta\in\{1,2\}\), and suppose \(b_1\ge2\). Assume that, for
every \(x\in X\), \(y\in Y\), and \(a\in\mathcal A\), the graph \(H\)
has an \(x,y\)-path of length \(a\). Assume also that, for every
\(b\in\mathcal B\), the graph \(G\) has an \((X,Y)\)-path of length
\(b\) whose internal vertices avoid \(V(H)\). Then \(G\) has at least
\(s+t-1\) admissible cycles.

The same conclusion holds with \(X=Y\) if every outside path has
distinct ends and, for every two distinct \(x,y\in X\) and every
\(a\in\mathcal A\), the required \(x,y\)-path of length \(a\) exists in
\(H\).

We also use a two-variable form. Let \(P_1,\ldots,P_s\) and
\(Q_1,\ldots,Q_t\) be path families with length sets \(\mathcal A\) and
\(\mathcal B\), respectively. Suppose that, for every \(i,j\), these
paths and fixed connectors form a simple cycle \(C_{ij}\), and that \[
 |E(C_{ij})|=|E(P_i)|+|E(Q_j)|+\lambda
\] for a constant \(\lambda\) independent of \(i,j\). Then the cycles
\(C_{ij}\) contain an admissible subfamily of at least \(s+t-1\)
members.

\Needspace{10\baselineskip}

\textbf{Proof.} For each \(b\in\mathcal B\), choose a corresponding
outside path \(P_b\), with ends \(x_b,y_b\). Its internal vertices avoid
\(V(H)\), whereas every core \(x_b,y_b\)-path lies in \(H\). The two
paths meet exactly at \(x_b,y_b\), and hence form a simple cycle whose
length is their total length. Thus every element of
\(\mathcal A+\mathcal B\) occurs as a cycle length, and Lemma 2.1 gives
the required admissible subfamily. The same argument applies when
\(X=Y\), because the outside path has distinct ends.

In the two-variable form, the cycle lengths contain
\(\lambda+\mathcal A+\mathcal B\), so Lemma 2.1 applies once more.
\(\square\)

In each use of the two-variable form, we identify the supports of the
variable paths and the fixed connectors. The stated intersection pattern
then proves both simplicity and the displayed length formula. In
particular, two three-term admissible families on such sides, with fixed
connectors, yield five admissible cycles by Lemma 2.1. Their lengths
include a multiple of five, since a five-term progression of difference
one or two contains every residue modulo five.

\textbf{Lemma 2.3 (the \(3+3\) lemma modulo five).} Let \(d\in\{1,2\}\),
let \(a\in\mathbb Z\), and let \(R\subseteq\mathbb Z/5\mathbb Z\) have
three distinct elements. Then \[
 \{a,a+d,a+2d\}+R
\] contains \(0\) modulo five.

\textbf{Proof.} Put \(R'=d^{-1}(a+R)\). Since \(d\in\{1,2\}\) is a unit
modulo five, \(R'\) has three elements. The two three-element subsets
\(R'\) and \(\{0,-1,-2\}\) of \(\mathbb Z/5\mathbb Z\) intersect. If
\(r'=-i\) lies in the intersection, with \(i\in\{0,1,2\}\), then
\(a+id+r=0\) for the corresponding \(r\in R\). \(\square\)

\textbf{Lemma 2.4 (theta residues in girth six).} Let \(H\) be a theta
in a graph of girth at least six. Then some two distinct vertices of
\(H\) are joined in \(H\) by three paths whose lengths are pairwise
distinct modulo five.

\textbf{Proof.} Let the three defining \(x,y\)-paths have lengths
\(a,b,c\). Each pair forms a cycle, so every pairwise sum is at least
six. If the three lengths are already distinct modulo five, there is
nothing to prove.

Suppose exactly two residues agree, say \(a\equiv b\not\equiv c\pmod5\),
and relabel so that \(a\ge3\). For \(r\in\{1,2\}\), let \(u_r\) be the
vertex at distance \(r\) from \(x\) on the first path. The three
\(u_r,y\)-paths that continue along the first path, or return to \(x\)
and then follow the second or third path, have lengths \[
 a-r,\qquad r+b,\qquad r+c \pmod5.
\] The first two residues are distinct, and so are the last two. The
first and third can agree for at most one of \(r=1,2\); the other value
of \(r\) therefore gives three distinct residues.

It remains that \(a\equiv b\equiv c\equiv t\pmod5\). Since every
pairwise sum is at least six, relabel so that \(a,b\ge3\). Choose
vertices at distances \(r\) and \(s\) from \(x\) on the first and second
paths. The routes through \(x\), through \(y\), and through the third
theta path have residues \[
 r+s,\qquad 2t-r-s,\qquad r+2t-s \pmod5.
\] For \(t=0,3,4\) take \((r,s)=(1,1)\); for \(t=1\) take \((1,2)\); and
for \(t=2\) take \((2,1)\). In each case the three residues are
distinct. All selected vertices are internal, and distinct theta paths
meet only at \(x,y\), so every route used above is simple. \(\square\)

\subsection{3. The rooted end-block
reduction}\label{the-rooted-end-block-reduction}

Assume for a contradiction that a nonempty finite simple graph \(G_0\)
has minimum degree at least five and no cycle of length divisible by
five. Choose a connected component \(G\) of \(G_0\). Every neighbor of a
vertex of \(G\) also lies in \(G\), so \(\delta(G)\ge5\); moreover,
every cycle in \(G\) is a cycle in \(G_0\).

Write \(n=|V(G)|\). If \(G\) is 2-connected, BGLP Theorem 1.3 with
\(k=5\) gives five admissible cycles, one of whose lengths is divisible
by five, unless \[
 G\cong K_6
 \qquad\text{or}\qquad
 G\cong K_{5,n-5}.
\] The first exception contains a 5-cycle. In the second, the
minimum-degree hypothesis forces both parts to have at least five
vertices, and hence gives a 10-cycle. Thus \(G\) is not 2-connected.

Choose an end block \(B\) of \(G\), with cut vertex \(c\). A leaf block
exists by the block-cut tree. It is not a bridge, since the non-cut
endpoint of an end bridge would have degree one in \(G\). Thus \(B\) is
2-connected, and \(c\) has at least two neighbors in \(B\). Every vertex
of \(B-c\) is inner in the end block, so all its incident edges in \(G\)
lie in \(B\). Therefore \[
 d_B(v)\ge5\qquad(v\in V(B)\setminus\{c\}). \tag{3.1}
\] If \(d_B(c)\ge5\), applying BGLP Theorem 1.3 to \(B\) again gives a
cycle whose length is divisible by five: one of the five admissible
cycles in the generic outcome has such a length, while the exceptional
graphs contain a 5-cycle or a 10-cycle. Hence \[
 2\le d_B(c)\le4. \tag{3.2}
\] Set \[
 J=B-c,
 \qquad
 D=\{v\in V(J):d_J(v)=4\}.
\] Deleting \(c\) removes at most one incident edge from each remaining
vertex. Thus \[
 d_J(v)=4\quad(v\in D),
 \qquad
 d_J(v)\ge5\quad(v\in V(J)\setminus D),
 \tag{3.3a}
\] and every vertex of \(D\) is adjacent to \(c\). In particular, \[
 D\subseteq N_B(c),
 \qquad
 |D|\le d_B(c)\le4. \tag{3.3b}
\] For the rest of the proof, \(B\) and \(J\) satisfy (3.1)--(3.3b).
Since \(B\subseteq G\subseteq G_0\), the graph \(B\) contains no cycle
of length divisible by five.

\subsubsection{3.1. Two-connectivity after deleting the
root}\label{two-connectivity-after-deleting-the-root}

\textbf{Lemma 3.1.} The graph \(J=B-c\) is 2-connected.

\textbf{Proof.} Since \(B\) is 2-connected, \(J\) is connected. Suppose
that \(J\) is not 2-connected. Choose distinct end blocks \(E_1,E_2\) of
\(J\), and let \(b_i\) be the cut vertex of \(J\) belonging to \(E_i\).
Neither end block is a bridge, because \(\delta(J)\ge4\).

For each \(i\), an inner vertex \(x_i\in V(E_i)\setminus\{b_i\}\) is
adjacent to \(c\); otherwise \(b_i\) would separate \(E_i-b_i\) from the
rest of \(B\). Every \(v\in V(E_i)\setminus\{x_i,b_i\}\) is inner in
\(E_i\), and hence \[
 d_{E_i}(v)=d_J(v)\ge4.
\] Apply GHLM Theorem 3.1 with \(q=3\) to \(E_i-x_ib_i\), rooted at
\(x_i,b_i\). The graph \((E_i-x_ib_i)+x_ib_i\) is 2-connected, and
deleting the edge \(x_ib_i\) changes only the degrees of the two roots.
We obtain three admissible \(x_i,b_i\)-paths with lengths \[
 \mathcal A_i=\{a_i,a_i+d_i,a_i+2d_i\},
 \qquad d_i\in\{1,2\}.
\]

The block-cut tree supplies a \(b_1,b_2\)-path \(R\) whose internal
vertices avoid the two end-block interiors; if \(b_1=b_2\), take \(R\)
to have length zero. Distinct end blocks have disjoint interiors and
meet, if at all, only in their common cut vertex. For paths \(P_i\)
chosen from the two families, \[
 c x_1+P_1+R+\overleftarrow{P_2}+x_2c
\] is a simple cycle: the two end-block interiors are disjoint, \(R\)
avoids both, and \(c\notin V(J)\). When \(b_1=b_2\), the two sides share
only \(c\) and their common endpoint. The possible lengths are \[
 2+|E(R)|+\mathcal A_1+\mathcal A_2.
\] By Lemma 2.1 these include five admissible cycle lengths, one
divisible by five. \(\square\)

\subsubsection{3.2. Three-connectivity}\label{three-connectivity}

\textbf{Lemma 3.2.} The graph \(J\) is 3-connected.

\textbf{Proof.} By Lemma 3.1, \(J\) is 2-connected, and (3.3a) gives
\(\delta(J)\ge4\). Since \(J\subseteq B\), it has no cycle whose length
is divisible by five and hence no five admissible cycles. BGLP Lemma 2.3
with \(k=5\) and \(r=1\) gives \[
 \kappa(J)\ge \frac{5+2}{2}-1=\frac52.
\] Since vertex connectivity is integral, \(\kappa(J)\ge3\). \(\square\)

\subsubsection{3.3. Root lifting}\label{root-lifting}

In every rooted auxiliary graph below, we shall verify the displayed
degree bounds. After its construction-specific boundary losses are
accounted for, only vertices of \(D\) may miss the required degree bound
by one. The next lemma records the repair: when there are at least two
such vertices, adjoining \(c\) restores every deficit and leaves \(c\)
as the single COY exception.

\textbf{Lemma 3.3 (root lifting).} Let \(A\) be a subgraph of \(J\),
with \(V(A)\subseteq V(J)\), and let \(x,y\in V(A)\) be distinct roots
such that \(xy\notin E(A)\) and \(A+xy\) is 2-connected. Let \(q\) be an
integer with \(2\le q\le4\), and let \[
 Z\subseteq D\cap\bigl(V(A)\setminus\{x,y\}\bigr)
\] satisfy \[
 d_A(v)\ge q+1
 \quad\bigl(v\in V(A)\setminus(\{x,y\}\cup Z)\bigr),
 \qquad
 d_A(z)\ge q\quad(z\in Z). \tag{3.4}
\] Assume additionally that \(|V(A)|\ge4\) in the case \(|Z|=1\). Then
\(B\) contains \(q\) admissible \(x,y\)-paths whose internal vertices
lie in \(V(A)\), except that when \(|Z|\ge2\) they may also contain
\(c\).

\textbf{Proof.} If \(Z=\varnothing\), apply GHLM Theorem 3.1 to \(A\).
If \(|Z|=1\), apply the COY theorem to \(A\), taking the unique vertex
of \(Z\) as the exception.

\Needspace{8\baselineskip}

Now suppose \(|Z|\ge2\). Let \(A^+\) be obtained from \(A\) by adjoining
\(c\) and the edges \(cz\) for \(z\in Z\). By (3.3b), each such edge
belongs to \(B\), so \(A^+\subseteq B\). The graph \(A^++xy\) is
2-connected: it is obtained from \(A+xy\) by adjoining a vertex with at
least two neighbors. Each vertex of \(Z\) gains one degree, every other
nonroot already has degree at least \(q+1\), and among the nonroots only
\(c\) remains exceptional. Since \(A+xy\) is 2-connected,
\(|V(A)|\ge3\), and hence \(|V(A^+)|\ge4\). The COY theorem gives the
required paths. \(\square\)

\subsection{4. From 3-connectivity to
4-connectivity}\label{from-3-connectivity-to-4-connectivity}

We next exclude a three-vertex cut. The first lemma supplies the rooted
2-connectivity on either side of such a cut.

\textbf{Lemma 4.1 (one component and two separator roots).} Let
\(S=\{x,y,z\}\) be a three-vertex cut of \(J\), let \(Q\) be a component
of \(J-S\), and put \[
 A(Q;x,y)=J[Q\cup\{x,y\}]-xy.
\] Then \(A(Q;x,y)+xy\) is 2-connected.

\textbf{Proof.} Each vertex of \(S\) has a neighbor in every component
of \(J-S\). By symmetry it suffices to consider \(z\). If \(z\) had no
neighbor in \(Q\), then \(J-\{x,y\}\) would separate \(Q\) from \(z\)
and the remainder of the graph, contrary to Lemma 3.2. Hence \(Q\)
together with either root is connected.

Let \(K=A(Q;x,y)+xy\). Since \(Q\ne\varnothing\) and \(x,y\notin Q\),
the graph \(K\) has at least three vertices; it is connected because
\(Q\) is connected and each root has a neighbor in \(Q\). Deleting
either root leaves the other joined to \(Q\). If some \(w\in Q\) were a
cut vertex of \(K\), the edge \(xy\) would keep the roots in one
component of \(K-w\), while another component \(R\subseteq Q-w\) would
have no neighbor in \(\{x,y\}\). Then \(J-\{w,z\}\) would separate \(R\)
from the rest of \(J\), again contradicting 3-connectivity. Thus \(K\)
is 2-connected. \(\square\)

Every component \(Q\) of \(J-S\) has at least two vertices, since a
singleton would have degree at most three. We shall also use the
following extension. If \(z\) has at least two neighbors in \(Q\), then
\[
 \bigl(J[Q\cup S]-xy\bigr)+xy
\] is 2-connected: start with \(A(Q;x,y)+xy\) and adjoin \(z\) through
two of its neighbors in \(Q\).

For later use, if \(Q\) is any component of \(J-S\) and \(A=A(Q;x,y)\),
then the only separator vertex omitted from \(A\) is \(z\). Deleting
\(xy\) does not affect a vertex of \(Q\), so \[
 d_A(v)\ge4\quad(v\in Q\setminus D),
 \qquad
 d_A(v)\ge3\quad(v\in Q\cap D). \tag{4.1}
\]

\textbf{Lemma 4.2.} If \(S\) is a three-vertex cut of \(J\), then every
component of \(J-S\) contains at least two vertices of \(D\).
Consequently, \(J-S\) has exactly two components, \(|D|=4\), and each
component contains exactly two vertices of \(D\).

\textbf{Proof.} Write \(S=\{x,y,z\}\). Suppose that a component \(Q_0\)
of \(J-S\) contains at most one vertex of \(D\), and let \(Q_1\) be
another component. For \(i\in\{0,1\}\), set \[
 A_i=A(Q_i;x,y).
\] \Needspace{8\baselineskip}

Let \(Z_i=D\cap Q_i\). Since \(|Q_i|\ge2\), Lemma 3.3 with \(q=3\)
applies and gives three admissible \(x,y\)-paths on each side: Lemma 4.1
gives the augmented 2-connectivity, (4.1) gives the degrees, and
\(|V(A_i)|=|Q_i|+2\ge4\) supplies the exceptional-case order. The
\(Q_0\)-family avoids \(c\), because \(|Z_0|\le1\); the \(Q_1\)-family
may use \(c\).

Apart from the roots, the two families lie in the disjoint components
\(Q_0,Q_1\), and at most one may contain \(c\). Any two selected paths
therefore form a simple cycle. Lemma 2.1 gives five admissible cycle
lengths, a contradiction. Thus every component of \(J-S\) contains at
least two vertices of \(D\). Since \(J-S\) has at least two components
and \(|D|\le4\), it has exactly two; each contains exactly two vertices
of \(D\), and \(D\cap S=\varnothing\). \(\square\)

\textbf{Lemma 4.3.} The graph \(J\) is 4-connected.

\textbf{Proof.} Suppose that \(S\) is a three-vertex cut. Lemma 4.2
shows that \(J-S\) has precisely two components \(Q_1,Q_2\) and gives \[
 |D\cap Q_1|=|D\cap Q_2|=2,
 \qquad
 D\cap S=\varnothing. \tag{4.2}
\] Every vertex of \(S\) therefore has degree at least five in \(J\).
Retain one separator vertex, say \(z\), and denote the other two by
\(x,y\). The vertex \(z\) has at most two neighbors in \(S\), so one of
\(Q_1,Q_2\) contains at least two of its neighbors; relabel that
component as \(Q_2\).

On the first side, put \[
 A_1=J[Q_1\cup\{x,y\}]-xy,
 \qquad Z_1=D\cap Q_1.
\] Lemma 4.1 and (4.1) allow Lemma 3.3 with \(q=3\). Since \(|Z_1|=2\),
it yields three admissible \(x,y\)-paths with interiors in
\(Q_1\cup\{c\}\).

On the second side, put \[
 A_2=J[Q_2\cup S]-xy.
\] By the extension following Lemma 4.1, \(A_2+xy\) is 2-connected.
Every vertex of \(Q_2\) retains all its \(J\)-neighbors in
\(Q_2\cup S\), and therefore has degree at least four in \(A_2\). The
only other vertices are the roots \(x,y\) and the exceptional vertex
\(z\). Since \(|V(A_2)|\ge5\), the COY theorem with \(q=3\) gives three
admissible \(x,y\)-paths with interiors in \(Q_2\cup\{z\}\).

The two families meet only at \(x,y\): their remaining vertices lie on
opposite sides of the cut, and only the first family can contain \(c\).
Pairwise unions are therefore simple cycles. Lemma 2.1 gives five
admissible lengths, a contradiction. Hence no three-cut exists. Since
\(\delta(J)\ge4\), the graph has at least five vertices, and Lemma 3.2
already excludes cuts of size at most two. Thus \(J\) is 4-connected.
\(\square\)

\subsection{\texorpdfstring{5. First local exclusions: \(K_4^-\) and
triangles}{5. First local exclusions: K\_4\^{}- and triangles}}\label{first-local-exclusions-k_4--and-triangles}

\subsubsection{\texorpdfstring{5.1. Excluding
\(K_4^-\)}{5.1. Excluding K\_4\^{}-}}\label{excluding-k_4-}

\textbf{Lemma 5.1.} The graph \(J\) contains no subgraph isomorphic to
\(K_4^-\).

\textbf{Proof.} Suppose that distinct vertices \(x,y,a,b\) contain the
five edges \[
 xy,\ xa,\ ya,\ yb,\ ab;
\] no condition is imposed on \(xb\). No vertex
\(w\in V(J)\setminus\{x,y,a,b\}\) is adjacent to both \(a\) and \(b\),
since then \[
 x-y-b-w-a-x
\] would be a 5-cycle. Put \(K=\{a,b\}\).

\Needspace{8\baselineskip}

Deleting two vertices from the 4-connected graph \(J\) leaves a
2-connected graph, so \(J-K\) is 2-connected. Let \[
 A=(J-K)-xy,
\] with roots \(x,y\). Adding \(xy\) gives a 2-connected graph. By the
preceding paragraph, every vertex of \(V(A)\setminus\{x,y\}\) has at
most one neighbor in \(K\), and hence \[
 d_A(v)\ge4\quad(v\in V(A)\setminus(D\cup\{x,y\})),
 \qquad
 d_A(v)\ge3\quad
   (v\in (D\cap V(A))\setminus\{x,y\}). \tag{5.1}
\] Set \(Z=D\cap(V(A)\setminus\{x,y\})\). If \(|V(A)|=3\), then
\(|V(J)|=5\), and 4-connectivity forces \(J=K_5\), which contains a
5-cycle. Since \(A+xy\) is 2-connected, \(|V(A)|\ge3\); the equality
case has just been excluded. Thus \(|V(A)|\ge4\), and Lemma 3.3 with
\(q=3\) gives three admissible \(x,y\)-paths whose interiors avoid
\(K\), apart from the possible vertex \(c\).

The subgraph on \(K\cup\{x,y\}\) contains \(x,y\)-paths of lengths
\(1,2,3\): \[
 xy,
 \qquad xay,
 \qquad xaby.
\] Each of these core paths meets every lifted path only at \(x,y\).
Lemma 2.2 therefore gives five admissible simple cycles in \(B\), a
contradiction. \(\square\)

\subsubsection{5.2. Excluding triangles}\label{excluding-triangles}

\textbf{Contraction fact.} Let \(G\) be a 4-connected simple graph and
let \(qr\in E(G)\). After contracting \(qr\) and suppressing loops and
parallel edges, the resulting simple graph \(G/qr\) is 3-connected.

\textbf{Proof.} Let \[
 \pi:V(G)\longrightarrow V(G/qr)
\] be the contraction map. Since \(G\) has at least five vertices,
\(G/qr\) has at least four. Let \(X\subseteq V(G/qr)\) with \(|X|\le2\),
and put \(\widetilde X=\pi^{-1}(X)\). If the contracted vertex is not in
\(X\), then \(|\widetilde X|=|X|\); if it is, then
\(|\widetilde X|=|X|+1\). Thus \(|\widetilde X|\le3\), and
\(G-\widetilde X\) is connected.

Take \(x,y\in V(G/qr)\setminus X\) and choose preimages
\(\widetilde x,\widetilde y\notin\widetilde X\). Apply \(\pi\) to the
vertex sequence of a \(\widetilde x,\widetilde y\)-path in
\(G-\widetilde X\) and delete consecutive repetitions. Consecutive
distinct images are adjacent in the simple quotient, so the resulting
sequence is an \(x,y\)-walk in \((G/qr)-X\). Every walk contains a path
with the same endpoints. Hence \((G/qr)-X\) is connected for every
\(|X|\le2\), and \(G/qr\) is 3-connected. \(\square\)

\textbf{Lemma 5.2.} The graph \(J\) is triangle-free.

\textbf{Proof.} Suppose that \(pqrp\) is a triangle. By Lemma 5.1, every
vertex of \(J-\{p,q,r\}\) has at most one neighbor in this triangle.

Choose an edge of the triangle, say \(qr\), whose contraction leaves
\(c\) with at least two distinct neighbors. Contracting one edge can
identify at most two neighbors of \(c\), so any edge works when
\(d_B(c)\ge3\). When \(d_B(c)=2\), only the edge joining those two
neighbors can identify them, and at most one triangle edge does so.

Contract \(qr\) to a vertex \(t\) in \(B\), suppress loops and parallel
edges, and delete \(pt\). Let the resulting graph be \(L\), rooted at
\(p,t\). The graph \(J/qr\) is 3-connected and hence 2-connected. By the
choice of \(qr\), the vertex \(c\) has two neighbors in \(J/qr\);
adjoining it through those neighbors shows that \(B/qr\) is 2-connected,
and \[
 L+pt=B/qr.
\] \Needspace{10\baselineskip}

For \(v\in V(B)\setminus\{p,q,r,c\}\), identify \(v\) with its unchanged
image in \(L\). The contraction merges no two edges incident with \(v\),
since \(v\) has at most one neighbor in the triangle. Hence \[
 d_L(v)=d_B(v)\ge5
 \qquad(v\in V(L)\setminus\{p,t,c\}).
\] Moreover \(|V(L)|\ge5\): the 4-connected graph \(J\) has at least
five vertices, \(B\) has the additional vertex \(c\), and the
contraction removes one vertex. The COY theorem with \(q=4\) and
exceptional vertex \(c\) gives four admissible \(p,t\)-paths in \(L\).

The edge \(pt\) is absent, so each rooted path has length at least two.
In a simple \(p,t\)-path, \(t\) occurs only at the end. Every other
quotient vertex has a unique unchanged preimage, and every preceding
path edge is an edge of \(B\). If \(v\) is the penultimate vertex, then
\(vt\) comes from \(vq\) or \(vr\) in \(B\); replace it by one such
edge. No earlier vertex represents \(q\) or \(r\), so this gives a
simple path of the same length ending at one of \(q,r\) and avoiding the
other.

Close a lift ending at \(q\) by \(qp\) or by \(qrp\), and use the
symmetric closures for a lift ending at \(r\). The unused triangle
vertex is absent from the lift, so each closing route meets it only at
its endpoints; both unions are simple cycles and add one or two edges.
If the four path lengths are \(a,a+d,a+2d,a+3d\), the closure lengths
are \[
 \{a+id+j:0\le i\le3,\ j\in\{1,2\}\}.
\] For \(d=1\) these contain \(a+1,\ldots,a+5\), and for \(d=2\) they
contain \(a+1,\ldots,a+8\). In either case one length is divisible by
five, a contradiction. \(\square\)

\subsection{6. Boundary lifting and the four-cycle
exclusion}\label{boundary-lifting-and-the-four-cycle-exclusion}

To exclude a four-cycle, we enlarge a complete bipartite core as far as
possible and construct admissible paths through a component outside it.
The next two lemmas justify the recurring construction: replace
attachments to the core by two artificial roots, apply a rooted path
theorem, and then replace the artificial boundary edges by genuine edges
back into the core.

\subsubsection{6.1. Boundary-root lifting}\label{boundary-root-lifting}

\textbf{Lemma 6.1.} Let \(H\) have distinct distinguished vertices
\(x,y\), and put \(Q=H-\{x,y\}\). Suppose that:

\begin{enumerate}
\def\labelenumi{\arabic{enumi}.}
\tightlist
\item
  \(Q\) is connected;
\item
  \(H\) contains independent edges \(xu\) and \(yv\) with distinct
  \(u,v\in V(Q)\); and
\item
  every end block of \(Q\) has an inner vertex adjacent in \(H\) to
  \(x\) or \(y\).
\end{enumerate}

Then \(H+xy\) is 2-connected.

\textbf{Proof.} The graph \(H+xy\) is connected, and deleting either
root leaves the other joined to \(Q\). The independent boundary edges
exhibit the four distinct vertices \(x,y,u,v\), so \(|V(H)|\ge4\) and
the order requirement for 2-connectivity holds. Suppose that \(w\in Q\)
is a cut vertex of \(H+xy\). The edge \(xy\) keeps the roots together,
and \(u\ne v\) ensures that at least one of their two attachments
survives the deletion of \(w\). Thus the roots lie with one component of
\(Q-w\), while some other component contains no neighbor of either root.
The block-cut tree supplies an end block of \(Q\) in that component and
away from \(w\), contradicting (3). \(\square\)

\Needspace{10\baselineskip}

\textbf{Lemma 6.2 (boundary-root lifting).} Let \(c\in V(B)\), put
\(J=B-c\), and let \(D\subseteq V(J)\) satisfy \(D\subseteq N_B(c)\).
Let \(H\) have distinct roots \(\alpha,\beta\), with
\(\alpha\beta\notin E(H)\) and \(c\notin V(H)\), and put \[
 Q=H-\{\alpha,\beta\}.
\] Suppose that \(Q\) admits an injective graph embedding into \(J\).
Relabel the old vertices by their images under this embedding;
henceforth \(V(Q)\subseteq V(J)\) and \(Q\subseteq J\). Suppose also
that \(H+\alpha\beta\) is 2-connected. Set \(Z=D\cap V(Q)\). Let \(q\)
be an integer with \(2\le q\le4\), and suppose \[
 d_H(v)\ge q+1\quad(v\in V(Q)\setminus Z),
 \qquad
 d_H(z)\ge q\quad(z\in Z). \tag{6.1}
\]

\Needspace{7\baselineskip}

Assume additionally that \(|V(H)|\ge4\) in the case \(|Z|=1\). Then:

\begin{enumerate}
\def\labelenumi{\arabic{enumi}.}
\tightlist
\item
  if \(|Z|\le1\), \(H\) contains \(q\) admissible \(\alpha,\beta\)-paths
  whose internal vertices lie in \(Q\);
\item
  if \(|Z|\ge2\), the graph \(H^+\) obtained by adjoining \(c\) and
  precisely the edges \(cz\), \(z\in Z\), contains \(q\) admissible
  \(\alpha,\beta\)-paths whose internal vertices lie in \(Q\cup\{c\}\).
\end{enumerate}

Suppose in addition that \(\alpha,\beta\notin V(B)\), and let
\(F\subseteq J\) satisfy \(V(F)\cap V(Q)=\varnothing\). For each edge
\(\alpha u\in E(H)\), choose a nonempty set \[
 L(u)\subseteq N_J(u)\cap V(F),
\] and for each edge \(v\beta\in E(H)\) choose a nonempty set \[
 R(v)\subseteq N_J(v)\cap V(F).
\] If, for every rooted path supplied above with first old vertex \(u\)
and last old vertex \(v\), one can choose \(x\in L(u)\) and
\(y\in R(v)\) with \(x\ne y\), then replacing \(\alpha u\) by \(xu\) and
\(v\beta\) by \(vy\) gives an \(x,y\)-path in \(B\) of the same length.
Its internal vertices lie in \(Q\cup\{c\}\) and avoid \(F\). The
endpoint pair may depend on the path, but the length sequence remains
admissible.

\textbf{Proof.} For \(Z=\varnothing\), apply GHLM Theorem 3.1 to \(H\).
For \(|Z|=1\), apply the COY theorem to \(H\), with the unique vertex of
\(Z\) as the exception. If \(|Z|\ge2\), then \(H^++\alpha\beta\) is
2-connected because it is obtained from \(H+\alpha\beta\) by adjoining
\(c\) through at least two neighbors. Each vertex of \(Z\) gains one
degree, every other nonroot already has degree at least \(q+1\), and
only \(c\) is exceptional. The two roots, two distinct vertices of
\(Z\), and \(c\) give \(|V(H^+)|\ge5\), so the COY order condition also
holds. The theorem gives the asserted paths.

For the boundary replacement, let \(P\) be one of these simple rooted
paths, and let \(u,v\) be the neighbors on \(P\) of \(\alpha,\beta\),
respectively. Deleting the artificial roots leaves a simple
\(u,v\)-path, reduced to the single vertex \(u\) when \(u=v\). Choose
distinct \(x\in L(u)\) and \(y\in R(v)\). Because \(F\) is disjoint from
\(Q\) and \(c\notin V(J)\), neither endpoint lies on the remaining
segment. Replacing \(\alpha u\) and \(v\beta\) by \(xu\) and \(vy\)
therefore preserves both simplicity and length. \(\square\)

\subsubsection{6.2. A lifted complete-bipartite
core}\label{a-lifted-complete-bipartite-core}

Let \(K_{s,t}^-\) denote a graph obtained from \(K_{s,t}\) by deleting
one edge, and set \[
 \mathcal F=
 \{K_{s,t}:\min\{s,t\}\ge2\}
 \cup
 \{K_{s,t}^-:\min\{s,t\}\ge2,\ \max\{s,t\}\ge3\}.
\]

\Needspace{11\baselineskip}

We shall use the following elementary path ranges.

\textbf{Complete-bipartite path lemma.} Let \(F=K_{s,t}\) with
\(r=\min\{s,t\}\ge2\), or let \(F=K_{s,t}^-\) with \(r\ge3\). Let
\((S,T)\) be its bipartition.

\begin{itemize}
\tightlist
\item
  If \(x,y\) lie in opposite parts and \(xy\in E(F)\), then \(F\) has an
  \(x,y\)-path of every odd length \(1,3,\ldots,2r-1\).
\item
  If \(x,y\) are distinct vertices in the same part, then \(F\) has an
  \(x,y\)-path of every even length \(2,4,\ldots,2r-2\). If that part
  has more than \(r\) vertices, it also has one of length \(2r\).
\end{itemize}

\textbf{Proof.} The complete case follows by choosing the required
number of vertices from each part and alternating. Suppose that the
missing edge is \(ab\), where \(a\in S\) and \(b\in T\); exchanging
\(S,T\) handles the reverse orientations below.

First let \(x,y\in S\) and let the desired length be \(2m\le2r-2\).
Choose distinct vertices \[
 s_0=x,s_1,\ldots,s_{m-1},s_m=y
 \quad\text{and}\quad
 t_0,\ldots,t_{m-1},
\] with every \(t_i\ne b\). This is possible because
\(|T\setminus\{b\}|\ge r-1\ge m\), and the alternating sequence \[
 s_0t_0s_1t_1\cdots s_{m-1}t_{m-1}s_m
\] is the required path.

For the terminal same-part length \(2r\), assume that the endpoint part,
say \(S\), has more than \(r\) vertices. Use an alternating sequence \[
 s_0t_0s_1t_1\cdots s_{r-1}t_{r-1}s_r,
 \qquad s_0=x,\quad s_r=y.
\] Avoid \(a\) or \(b\), the endpoints of the deleted edge, whenever the
vertex in question is not one of the prescribed path endpoints and
enough vertices remain in its part. If both \(a\) and \(b\) must occur,
place \(b\) at \(t_1\) when \(a=x\), at \(t_0\) when \(a=y\), and
otherwise place \(a\) at \(s_1\) and \(b\) at \(t_2\). Since \(r\ge3\),
these positions exist, and in each case \(a\) and \(b\) are
nonconsecutive.

Finally let \(x\in S\), \(y\in T\), and let the desired length be
\(2m+1\). The case \(m=0\) is the assumed edge \(xy\). Otherwise choose
an alternating sequence \[
 s_0t_0s_1t_1\cdots s_mt_m,
 \qquad s_0=x,\quad t_m=y.
\] Whenever possible, omit an unprescribed one of \(a,b\). If both are
forced to occur, then \(m\ge2\) and they can be placed as follows: if
\(a=x\), put \(b=t_1\); if \(b=y\), put \(a=s_1\); and if neither is
prescribed, put \(a=s_m\) and \(b=t_0\). The endpoints cannot satisfy
both \(a=x\) and \(b=y\) because \(xy\) is present. Again the missing
pair is nonconsecutive, so every displayed sequence is a path in \(F\).
\(\square\)

\textbf{Lemma 6.3 (lifted core lemma).} Under the standing setup
\(J=B-c\) and \(D=\{v\in V(J):d_J(v)=4\}\), assume \[
 \kappa(J)\ge4,
 \qquad J\text{ is triangle-free},
 \qquad \delta(J)\ge4,
 \qquad D\subseteq N_B(c),
 \qquad |D|\le4.
\] Suppose that \(F\subseteq J\) is isomorphic to a graph in
\(\mathcal F\), with bipartition \((S,T)\), and put
\(r=\min\{|S|,|T|\}\). If every vertex of \(J-V(F)\) has at most \(r-1\)
neighbors in \(F\), then \(B\) contains a cycle of length divisible by
five.

\textbf{Proof.} Suppose that \(B\) has no cycle of length divisible by
five. Let \(\mathcal C\) be the nonempty family of cores
\(F'\in\mathcal F\) for which every exterior vertex has at most
\(r(F')-1\) neighbors in \(F'\), where \(r(F')\) is its smaller part
size. Choose \(F\in\mathcal C\) of minimum order, label its parts
\((S,T)\), and reset \(r=|S|\le|T|\).

Triangle-freeness forbids edges within either part: for any same-part
edge, an opposite-part vertex not involved in the possible missing edge
completes a triangle. Thus the induced graph on the same vertex set is
either \(F\) or the complete bipartite graph obtained by restoring the
one missing edge. Replacing \(F\) by that induced graph preserves both
its order and the outside-neighbor bound. We may therefore assume that
\(F\) is induced.

If \(r\ge5\), then \(F\) contains a 10-cycle. This is immediate for a
complete core. If the missing edge is \(s_1t_1\), choose five vertices
in each part and use \[
 s_1t_2s_2t_1s_3t_3s_4t_4s_5t_5s_1.
\] Hence \[
 2\le r\le4. \tag{6.2}
\] The core is not spanning. If \(J=F=K_{s,t}^-\), then
\(\delta(J)\ge4\) forces \(r\ge5\). If \(J=F=K_{s,t}\), then (6.2) and
\(\delta(J)\ge4\) force \(r=4\); but then more than four vertices have
degree four, contrary to \(|D|\le4\).

Fix a component \(Q\) of \(J-V(F)\). It has at least two vertices, since
a singleton would have degree at most \(r-1\le3\). Every end block of
\(Q\) has an inner vertex adjacent to \(F\). If \(Q\) has no cut vertex
(including the case \(Q\cong K_2\)), then \(Q\) is its unique end block,
and the assertion follows from the connectedness of \(J\). Otherwise, if
an end block \(E\) had no such inner attachment and \(b\) were its
unique cut vertex, then \(J-b\) would separate \(V(E)\setminus\{b\}\)
from the rest of \(J\), contradicting 4-connectivity. This supplies the
end-block condition in Lemma 6.1 whenever all core attachments are
represented; the one construction that omits a boundary vertex is
checked separately below. Every auxiliary graph has at least four
vertices, so the order condition in Lemma 6.2 causes no further issue.

In every case below, an auxiliary graph retains all edges of \(Q\),
deletes the boundary edges from \(Q\) to the core, and replaces all
boundary edges of each represented type at a vertex by one edge to the
corresponding artificial root.

Each case now follows the same pattern. Core attachments are represented
by two artificial roots; the outside-neighbor bound gives the degree
estimate; Lemma 6.2 returns admissible paths to genuine, distinct
endpoints of \(F\); and the bipartite path ranges close them through the
core. It remains to distinguish whether \(Q\) has attachments in one or
both parts.

The numerical bookkeeping is the same in every case. The middle columns
count available lengths, and the last records the admissible-term
guarantee from Lemma 2.1: \[
\setlength{\arraycolsep}{3pt}
\begin{array}{c|c|c|c}
\text{case}
&\text{outside count}
&\text{core count}
&\text{guarantee}\\ \hline
r=2,\ Q\text{ has attachments in both parts}
&4&2&4+2-1=5\\
r=2,\ Q\text{ has attachments in one part}
&4&2&4+2-1=5\\
3\le r\le4,\ Q\text{ has attachments in both parts}
&6-r&r&(6-r)+r-1=5\\
3\le r\le4,\ Q\text{ has attachments in one part}
&7-r&r-1&(7-r)+(r-1)-1=5
\end{array}
\]

\Needspace{10\baselineskip}

\paragraph{Rank two}\label{rank-two}

Suppose \(r=2\). If \(F=K_{2,t}^-\), let \(x_1y_1\) be its missing edge,
with \(x_1\in S\) and \(y_1\in T\). Deleting \(y_1\) leaves a smaller
\(K_{2,t-1}\) in \(\mathcal F\), since \(t\ge3\), with the same
outside-neighbor bound: the newly external vertex \(y_1\) has exactly
one neighbor in the smaller core, and shrinking the core cannot increase
any other exterior count. This contradicts the minimality of \(F\). Thus
\(F=K_{2,t}\).\footnote{This is also the rank-two reduction used in the
  complete-bipartite step of the proof of BGLP Theorem 1.3; every
  subsequent deleted-edge use in that proof has smaller part at least
  three.}

\Needspace{10\baselineskip}

Since \(J\) is 3-connected and \(|S|=2\), the component \(Q\) must have
an attachment in \(T\). Suppose that it also has an attachment in \(S\).
Form \(H_{ST}\) from \(Q\) by adjoining roots \(s^*,t^*\) and joining
\(v\in Q\) to a root whenever \(v\) has a neighbor in the represented
part. The outside-neighbor bound is one, so the two root attachments
have distinct endpoints in \(Q\). Lemma 6.1 gives that \(H_{ST}+s^*t^*\)
is 2-connected. Each boundary edge deleted at a vertex of \(Q\) is
replaced by the corresponding root edge, and hence \[
 d_{H_{ST}}(v)=d_J(v)\ge5\quad(v\in V(Q)\setminus D),
 \qquad
 d_{H_{ST}}(v)=4\quad(v\in V(Q)\cap D).
\] \Needspace{8\baselineskip}

For an edge \(s^*v\), use \(N_J(v)\cap S\) as its represented endpoint
set; for an edge \(vt^*\), use \(N_J(v)\cap T\). The represented
endpoints lie in opposite parts and are therefore distinct. Lemma 6.2
with \(q=4\) gives four admissible outside paths, while the core gives
lengths one and three between every corresponding pair of endpoints.
Lemma 2.2 yields five admissible cycles.

We may therefore assume that every attachment of \(Q\) lies in \(T\).
Three-connectivity gives at least three attachment vertices in \(T\), so
\(|T|\ge3\). The bipartite attachment graph between \(Q\) and \(T\) has
a matching of size two. Otherwise a bipartite graph with no two
independent edges is a star, so all attachment edges share one endpoint.
If that endpoint lies in \(T\), deleting it separates \(Q\); if it is a
vertex \(v\in Q\), deleting \(v\) separates the nonempty set \(Q-\{v\}\)
from the core. Choose independent attachment edges with distinct ends
\(y_1,y_2\in T\), and partition \(T=T_1\cup T_2\) with \(y_i\in T_i\).

Form \(H_{12}\) by adjoining roots \(t_1^*,t_2^*\) and joining
\(v\in Q\) to \(t_i^*\) when \(v\) has a neighbor in \(T_i\). Lemma 6.1
applies, and the same degree calculation gives \(d_{H_{12}}(v)=d_J(v)\).
Represent the two boundary edges by the corresponding neighborhoods in
\(T_1\) and \(T_2\). Lemma 6.2 with \(q=4\) gives four admissible
outside paths. Between any endpoint in \(T_1\) and any endpoint in
\(T_2\), the core \(K_{2,t}\) has paths of lengths two and four; the
latter uses a third vertex of \(T\). Lemma 2.2 again gives five
admissible cycles. Both rank-two cases are impossible, so \[
 r\ge3. \tag{6.3}
\]

\paragraph{Attachments in both parts}\label{attachments-in-both-parts}

Suppose that \(Q\) has attachments in both \(S\) and \(T\). First let
\(F=K_{s,t}\). Form \(H_{\mathrm{opp}}\) by adjoining roots that
represent \(S\) and \(T\). No vertex of \(Q\) has neighbors in both
parts, since any such pair of neighbors would form a triangle with their
cross-edge. Thus the root attachments are independent, and Lemma 6.1
applies. For \(v\in Q\), put \(m=|N_J(v)\cap V(F)|\le r-1\). If \(m>0\),
replacing the \(m\) core edges by one root edge gives degree
\(d_J(v)-m+1\); if \(m=0\), the degree is unchanged. Hence \[
 d_{H_{\mathrm{opp}}}(v)\ge d_J(v)-r+2
 \ge
 \begin{cases}
  7-r,&v\in V(Q)\setminus D,\\
  6-r,&v\in V(Q)\cap D.
 \end{cases} \tag{6.4}
\] Use the original neighborhoods in \(S\) and \(T\) as the represented
endpoint sets. Lemma 6.2 with \[
 q=6-r\in\{2,3\}
\] gives \(6-r\) admissible outside paths between the two parts. The
core supplies the \(r\) odd lengths \(1,3,\ldots,2r-1\) between every
endpoint pair. Since \((6-r)+r-1=5\), Lemma 2.2 gives a contradiction.

Now let \(F=K_{s,t}^-\), with missing edge \(x_1y_1\), where
\(x_1\in S\) and \(y_1\in T\). Some core endpoint of an attachment edge
from \(Q\) lies outside \(\{x_1,y_1\}\). Otherwise \(J-\{x_1,y_1\}\)
would separate \(Q\) from the nonempty graph \(F-\{x_1,y_1\}\).
Interchanging \(S,T\) if necessary, assume that \(Q\) has a neighbor in
\(T\setminus\{y_1\}\); by the present case assumption, it also has a
neighbor in \(S\). Form \(H_{\mathrm{del}}\) with roots representing
\(S\) and \(T\setminus\{y_1\}\).

\Needspace{8\baselineskip}

Every end block \(E\) of \(Q\) has an inner vertex adjacent to
\(F-\{y_1\}\). If \(Q\) has no cut vertex (including \(Q\cong K_2\)),
then it is its unique end block, and the chosen attachment in
\(T\setminus\{y_1\}\) supplies such a vertex. Otherwise, let \(b\) be
the unique cut vertex of an end block \(E\). If \(E-b\) had no neighbor
in \(F-\{y_1\}\), then no edge of \(J-\{b,y_1\}\) would leave \(E-b\):
inside \(Q\), every such edge passes through \(b\), while by assumption
its only possible core attachment is \(y_1\). Thus \(\{b,y_1\}\) would
separate \(J\), contrary to 4-connectivity.

Finally, no vertex of \(Q\) is adjacent both to a vertex of \(S\) and to
a vertex of \(T\setminus\{y_1\}\), since those two core vertices are
adjacent and would form a triangle with it. Hence the two root
attachments have distinct ends in \(Q\). Lemma 6.1 applies.

For \(v\in Q\), put \(m=|N_J(v)\cap V(F)|\le r-1\). If \(v\) has a
represented neighbor, the new graph adds a root edge and \[
 d_{H_{\mathrm{del}}}(v)\ge d_J(v)-m+1\ge d_J(v)-r+2.
\] If it has none, then all its core neighbors lie in \(\{y_1\}\), so
\(d_{H_{\mathrm{del}}}(v)\ge d_J(v)-1\ge d_J(v)-r+2\). Thus the bounds
in (6.4) also hold for \(H_{\mathrm{del}}\). Lemma 6.2 supplies \(6-r\)
admissible paths from \(S\) to \(T\setminus\{y_1\}\). Every resulting
endpoint pair is joined by an edge of \(F\), and the complete-bipartite
path lemma applies because \(r\ge3\); it supplies the \(r\) odd lengths
\[
 1,3,\ldots,2r-1.
\] Since \((6-r)+r-1=5\), Lemma 2.2 gives five admissible cycles. We
conclude that \(Q\) has neighbors in only one part of \(F\).

\paragraph{Attachments in one part}\label{attachments-in-one-part}

By symmetry, assume that every attachment of \(Q\) lies in \(T\). The
attachment graph between \(Q\) and \(T\) has a matching of size two;
otherwise it is a star. Deleting its center disconnects all of \(Q\)
when the center lies in \(T\), or the nonempty set \(Q-\{v\}\) when the
center is \(v\in Q\). Choose independent attachment edges with distinct
ends \(y_1,y_2\in T\), and partition \(T=T_1\cup T_2\) with
\(y_i\in T_i\).

Construct \(H_{\mathrm{same}}\) from \(Q\) by adjoining roots
\(t_1^*,t_2^*\). For \(v\in Q\), add \[
 vt_1^*\quad\text{if }d_{T_1}(v)\ge1\text{ or }d_T(v)\ge2,
\] and \[
 vt_2^*\quad\text{if }d_{T_2}(v)\ge1\text{ or }d_T(v)\ge2.
\] The chosen matching gives independent root attachments, and every
end-block attachment is represented; hence Lemma 6.1 gives that
\(H_{\mathrm{same}}+t_1^*t_2^*\) is 2-connected. If \(d_T(v)\le1\),
every deleted boundary edge is replaced. If \(d_T(v)\ge2\), two root
edges are added and \(d_T(v)\le r-1\). Hence \[
 d_{H_{\mathrm{same}}}(v)\ge d_J(v)-r+3
 \ge
 \begin{cases}
  8-r,&v\in V(Q)\setminus D,\\
  7-r,&v\in V(Q)\cap D.
 \end{cases} \tag{6.5}
\]

For each root edge \(vt_i^*\), define its represented endpoint set by \[
 L_i(v)=
 \begin{cases}
 N_T(v),&d_T(v)\ge2,\\[2mm]
 N_T(v)\cap T_i,&d_T(v)=1.
 \end{cases}
\] \Needspace{8\baselineskip}

This set is nonempty whenever \(vt_i^*\) is present. We claim that, for
every rooted \(t_1^*,t_2^*\)-path, the two represented endpoints can be
chosen distinct.

Let \(u\) and \(v\) be the first and last old vertices of the rooted
path. If \(u=v\), then the presence of both root edges forces
\(d_T(u)\ge2\): a unique \(T\)-neighbor belongs to only one of
\(T_1,T_2\). We may therefore choose two distinct neighbors of \(u\). If
\(u\ne v\) and both represented sets are singletons, their elements lie
in the disjoint sets \(T_1,T_2\). If at least one represented set has at
least two elements, choose the endpoint from the other set first and
then avoid it in the larger set. This proves the boundary-replacement
condition in Lemma 6.2.

\Needspace{13\baselineskip}

Lemma 6.2 with \[
 q=7-r\in\{3,4\}
\] therefore gives \(7-r\) admissible paths with distinct ends in \(T\)
and interiors in \(Q\cup\{c\}\). Between any two distinct vertices of
\(T\), the core supplies the \(r-1\) even lengths \[
 2,4,\ldots,2r-2.
\] Since \((7-r)+(r-1)-1=5\), Lemma 2.2 gives five admissible cycles.
This final case is also impossible, proving the lemma. \(\square\)

\subsubsection{6.3. Excluding four-cycles}\label{excluding-four-cycles}

\textbf{Lemma 6.4.} The graph \(J\) contains no four-cycle.

\textbf{Proof.} A four-cycle contains a \(K_{2,2}\). Let \(s\ge2\) be
maximal such that \(K_{s,s}\subseteq J\). Since \(K_{5,5}\) contains a
10-cycle, \(s\le4\).

Suppose first that \(J\) contains \(K_{s+1,s+1}^-\) with bipartition
\((S,T)\). Maximality of \(s\) excludes \(K_{s+1,s+1}\), so its
designated missing edge \(x_0y_0\) is genuinely absent. An outside
vertex has at most \(s\) neighbors in either part: a vertex adjacent to
all of \(S\), for example, could replace \(y_0\) and complete a
\(K_{s+1,s+1}\). Moreover, a vertex meeting both parts can meet only
\(x_0\) and \(y_0\), since any other cross-pair is an edge of the core
and would complete a triangle. Thus every outside vertex has at most
\(s\) neighbors in the whole core. Lemma 6.3, with \(r=s+1\), gives a
contradiction. Hence \(J\) contains no \(K_{s+1,s+1}^-\).

Suppose next that \(K_{s,s+1}\subseteq J\). Fix its \(s\)-vertex side
\(S\), and choose a largest set \(T\subseteq V(J)\setminus S\) such that
all edges between \(S\) and \(T\) are present. Put \(t=|T|\), and let
\(F\) be the resulting \(K_{s,t}\). Triangle-freeness prevents an
outside vertex from meeting both parts. Maximality of \(t\) prevents it
from meeting all of \(S\). If an outside vertex \(v\) had \(s\)
neighbors \(T_0\subseteq T\), choose \(y\in T\setminus T_0\), possible
because \(t\ge s+1\). Then the bipartition \[
 \bigl(S\cup\{v\},\ T_0\cup\{y\}\bigr)
\] has every required cross-edge except possibly \(vy\). If \(vy\) is
absent, these vertices contain a \(K_{s+1,s+1}^-\), already excluded; if
it is present, they contain a \(K_{s+1,s+1}\), contradicting the
maximality of \(s\). Thus every outside vertex has at most \(s-1\)
neighbors in \(F\), and Lemma 6.3 gives a contradiction. Hence \(J\)
contains no \(K_{s,s+1}\).

Finally, take \(F=K_{s,s}\). An outside vertex meets at most one part
and cannot meet all \(s\) vertices there, since that would create a
\(K_{s,s+1}\). Lemma 6.3 applies once more and gives the final
contradiction. \(\square\)

The structural reduction is now complete. Lemmas 5.2 and 6.4 exclude
cycles of lengths three and four, while the counterexample assumption
excludes 5- and 10-cycles in \(J\subseteq B\). Therefore \[
 \operatorname{girth}(J)\ge6,
 \qquad
 J\text{ has no 10-cycle}. \tag{6.6}
\]

\subsection{7. The minimum-theta
argument}\label{the-minimum-theta-argument}

We now use (6.6) for the final contradiction. Since \(J\) is 2-connected
and \(\delta(J)\ge4\), it is not a cycle and therefore contains a theta.
Choose one, \(H\), of minimum order. By GHLM Lemma 5.10, \(H\) is
induced, at most one vertex outside \(H\) has two or more neighbors in
\(H\), and adjoining that vertex, if it exists, produces an induced
one-subdivision of \(K_4\).

\textbf{Subdivision attachment lemma.} Let \(K\) be a one-subdivision of
\(K_4\) in a graph \(A\) of girth at least five. If
\(v\in V(A)\setminus V(K)\) has two distinct neighbors in \(K\), then
\(A\) contains a cycle of length five or ten.

\textbf{Proof.} Write \(a_1,a_2,a_3,a_4\) for the branch vertices of
\(K\), and let \(z_{ij}=z_{ji}\) be the subdivision vertex on the edge
\(a_ia_j\) of the underlying \(K_4\). Let \(x,y\) be two distinct
neighbors of \(v\) in \(K\).

If both are branch vertices, say \(x=a_i\) and \(y=a_j\), then \[
 v a_i z_{ij} a_j v
\] is a 4-cycle, contrary to the girth assumption.

Suppose next that \(x=a_i\) is a branch vertex and \(y=z_{jk}\) is a
subdivision vertex. If \(i\in\{j,k\}\), then \(x\) and \(y\) are
adjacent and form a triangle with \(v\). Otherwise \(i,j,k\) are
distinct, and \[
 v a_i z_{ij} a_j z_{jk} v
\] is a 5-cycle.

It remains that \(x=z_{ij}\) and \(y=z_{k\ell}\) are subdivision
vertices. If the corresponding edges of \(K_4\) share an endpoint, then
\(x\) and \(y\) are joined by a two-edge path through that branch
vertex, giving a 4-cycle with \(v\). Otherwise the two edges are
disjoint. Relabeling, take \(x=z_{12}\) and \(y=z_{34}\). Then \[
 z_{12}a_1z_{13}a_3z_{23}a_2z_{24}a_4z_{34}
\] is an eight-edge path in \(K\), and adjoining the two edges from its
ends to \(v\) gives a 10-cycle. \(\square\)

Define \(T=H\) if no such vertex exists; otherwise let \(T\) be the
induced one-subdivision of \(K_4\) obtained by adjoining the unique such
vertex. For \(Q\subseteq V(J)\setminus V(T)\), write \[
 N_T(Q)=N_J(Q)\cap V(T).
\] If no vertex outside \(H\) has two neighbors in \(H\), then \(T=H\),
and GHLM Lemma 5.10 already gives \[
 |N_T(v)|\le1\qquad(v\in V(J)\setminus V(T)).
\] Otherwise \(T\) is a one-subdivision of \(K_4\). In that case, a
vertex outside \(T\) with two neighbors in \(T\) would, by the
subdivision attachment lemma, create a 5-cycle or a 10-cycle, contrary
to (6.6). Consequently, \[
 T\text{ is induced and 2-connected},
 \qquad
 \Delta(T)\le3,
 \qquad
 |N_T(v)|\le1\quad(v\in V(J)\setminus V(T)). \tag{7.1}
\] By Lemma 2.4, there are distinct \(x,y\in V(H)\subseteq V(T)\) joined
in \(T\) by three paths whose lengths are pairwise distinct modulo five.

Set \(M=J-V(T)\). The graph \(M\) is nonempty, since otherwise \(J=T\)
would have maximum degree at most three. By (3.3a) and (7.1), each
vertex loses at most one neighbor when \(T\) is deleted, so \[
 d_M(v)\ge4\quad(v\in V(M)\setminus D),
 \qquad
 d_M(v)\ge3\quad(v\in V(M)\cap D). \tag{7.2}
\] We show first that every component of \(M\) is 2-connected and then
that \(M\) has only one component.

\subsubsection{7.1. The components of the
exterior}\label{the-components-of-the-exterior}

\textbf{Lemma 7.1.} Every component of \(M\) is 2-connected.

\textbf{Proof.} Suppose that a component \(Q\) is not 2-connected. Since
\(Q\) is a component of \(M\), (7.2) gives minimum degree at least three
in \(Q\); hence every end block is nontrivial. Choose distinct end
blocks \(E_1,E_2\), and let \(b_i\) be the cut vertex of \(Q\) belonging
to \(E_i\).

For each \(i\), \[
 |N_T(E_i-b_i)|\ge3. \tag{7.3}
\] Indeed, if this set had at most two vertices, deleting it together
with \(b_i\) would separate the nonempty interior \(E_i-b_i\) from the
rest of \(J\). No edge leaves that interior inside \(Q\) except through
\(b_i\), and at least one vertex of the 2-connected graph \(T\)
survives. This contradicts the 4-connectivity of \(J\).

Choose attachment edges \(u_iy_i\) with \(u_i\in E_i-b_i\) and
\(y_i\in T\), with \(y_1\ne y_2\); the three choices available in (7.3)
permit this. Distinct end blocks have disjoint interiors, although their
cut vertices may coincide. The block-cut tree gives a \(b_1,b_2\)-path
\[
 R\subseteq Q-\bigl((E_1-b_1)\cup(E_2-b_2)\bigr);
\] take \(R\) to be the length-zero path when \(b_1=b_2\). Set \[
 A_i=E_i-u_ib_i,
 \qquad
 Z_i=D\cap\bigl(V(E_i)\setminus\{u_i,b_i\}\bigr).
\] Call \(E_i\) \textbf{light} if \(|Z_i|\le1\) and \textbf{heavy}
otherwise.

Adding the root edge in \(A_i\) gives a 2-connected graph containing
\(E_i\). An inner vertex of an end block has no neighbor in \(Q\)
outside that block. Hence, for every \(v\in V(E_i)\setminus\{b_i\}\), \[
 d_{E_i}(v)=d_Q(v)=d_M(v).
\] Passing from \(E_i\) to \(A_i\) by deleting \(u_ib_i\) affects only
\(u_i,b_i\). Thus (7.2) gives the degree hypotheses of Lemma 3.3 with
\(q=3\). If \(|Z_i|=1\), its order condition also holds:
\(d_{E_i}(u_i)=d_M(u_i)\ge3\) by (7.2), so \(|V(A_i)|=|V(E_i)|\ge4\).
Thus each \(E_i\) supplies three admissible \(u_i,b_i\)-paths. A light
family avoids \(c\), whereas a heavy family may use \(c\).

\Needspace{12\baselineskip}

Suppose first that at least one end block is light. Then, for every
choice of paths \(P_i\) from the two admissible families, at most one of
\(P_1,P_2\) contains \(c\). Fix a \(y_1,y_2\)-path \(L\) in \(T\),
orient \(R\) from \(b_1\) to \(b_2\) and \(L\) from \(y_1\) to \(y_2\),
and orient each \(P_i\) from \(u_i\) to \(b_i\). Moreover, \[
 V(P_i)\subseteq V(E_i)\cup\{c\},
 \qquad
 V(R)\subseteq
 V(Q)\setminus\bigl((E_1-b_1)\cup(E_2-b_2)\bigr),
 \qquad
 V(L)\subseteq V(T).
\] The interiors of \(E_1,E_2\) are disjoint; \(R\) meets \(P_i\) only
at \(b_i\); \(L\) is disjoint from \(M\cup\{c\}\); and \(y_1\ne y_2\).
If \(b_1=b_2\), that common cut vertex is the prescribed junction
through the length-zero path \(R\). Thus \[
 P_1+R+\overleftarrow{P_2}+u_2y_2+\overleftarrow L+y_1u_1
\] is a simple cycle. If the selected path lengths are \(p_1,p_2\), its
length is \[
 p_1+p_2+|E(R)|+|E(L)|+2.
\] The final three terms are fixed, so the two-variable form of Lemma
2.2 gives five admissible cycles, a contradiction.

It remains that both end blocks are heavy. Let \(K_i\) be obtained from
\(E_i\) by adjoining \(c\) and all edges from \(c\) to \(D\cap V(E_i)\).
Since \(Z_i\) contains at least two vertices of \(E_i-b_i\), \(c\) has
at least two neighbors in \(E_i\), and adjoining it through those
distinct neighbors makes \(K_i\) 2-connected. Apply GHLM Theorem 3.1
with \(q=3\) to \(K_i-cb_i\), rooted at the distinct vertices \(c,b_i\).
Adding the root edge preserves 2-connectivity, and deleting it affects
only the roots. Every \(v\in V(E_i)\setminus(D\cup\{b_i\})\) has degree
at least four in \(E_i\), while every
\(v\in(D\cap V(E_i))\setminus\{b_i\}\) has degree at least three and
gains its edge to \(c\). Hence each side supplies three admissible
\(c,b_i\)-paths.

For each \(i\), the selected \(c,b_i\)-paths lie in \(E_i\cup\{c\}\).
Two such paths, one from each end block, have no common internal vertex:
their \(E_i-b_i\) portions lie in disjoint end-block interiors, and
their only common vertex is \(c\), together possibly with the common
endpoint \(b_1=b_2\). The path \(R\) meets them only at \(b_1,b_2\).
Consequently their union with \(R\) is a simple cycle.

The edge \(cb_i\) was deleted before applying the rooted theorem, so
each selected path has length at least two. If their lengths are
\(p_1,p_2\), the resulting cycle has length \[
 p_1+p_2+|E(R)|.
\] The two-variable form of Lemma 2.2 again produces five admissible
cycles, a contradiction. Hence every component of \(M\) is 2-connected.
\(\square\)

\Needspace{20\baselineskip}

\subsubsection{7.2. Connectivity of the
exterior}\label{connectivity-of-the-exterior}

\textbf{Two-link lemma.} Let \(T\) be 2-connected, and let
\(X,Y\subseteq V(T)\) be disjoint two-element sets. Then \(T\) contains
two vertex-disjoint \((X,Y)\)-paths whose ends are distinct in \(X\) and
distinct in \(Y\), and whose internal vertices lie outside \(X\cup Y\).

\textbf{Proof.} Add a new vertex \(a\) adjacent to the two vertices of
\(X\), and then a new vertex \(b\) adjacent to the two vertices of
\(Y\). Each new vertex is adjoined through two distinct neighbors, so
the resulting graph remains 2-connected. By Menger's theorem, it
contains two internally vertex-disjoint \(a,b\)-paths.

Delete \(a,b\). The two remaining paths are vertex-disjoint. Their first
vertices are distinct and therefore exhaust \(X\), since a common first
vertex would have been internal to both original paths. Similarly, their
last vertices are distinct and exhaust \(Y\). Neither remaining path can
contain another vertex of \(X\cup Y\) internally: that vertex is an
endpoint of the other path, contradicting vertex-disjointness. These are
the required links. \(\square\)

\textbf{Lemma 7.2.} The graph \(M\) is connected.

\textbf{Proof.} Suppose that \(Q_1,Q_2\) are distinct components of
\(M\). Each has at least four neighbors in \(T\), since deleting at most
three such neighbors would separate that component from the other. By
(7.1), no vertex of \(Q_i\) is incident with two attachment edges to
\(T\). Choose distinct \(y_{1,1},y_{1,2}\in N_T(Q_1)\) and then distinct
\[
 y_{2,1},y_{2,2}\in
 N_T(Q_2)\setminus\{y_{1,1},y_{1,2}\};
\] the latter choice is possible because \(|N_T(Q_2)|\ge4\). Choose
corresponding attachment edges \[
 u_{i,1}y_{i,1},
 \qquad
 u_{i,2}y_{i,2}
\] for \(i=1,2\). By construction the four interface vertices
\(y_{i,j}\) are pairwise distinct. By (7.1), \(u_{i,1}\ne u_{i,2}\) for
each \(i\), so both rooted endpoint pairs are distinct.

By Lemma 7.1, each \(Q_i\) is 2-connected. Set \[
 A_i=Q_i-u_{i,1}u_{i,2},
 \qquad
 Z_i=D\cap\bigl(V(Q_i)\setminus\{u_{i,1},u_{i,2}\}\bigr),
\] and call \(Q_i\) light if \(|Z_i|\le1\) and heavy otherwise. Adding
the root edge in \(A_i\) gives a 2-connected graph containing \(Q_i\).
Since \(Q_i\) is a component of \(M\), every vertex retains all its
\(M\)-neighbors, and deleting the root edge affects only the roots; thus
(7.2) gives the degree hypotheses of Lemma 3.3 with \(q=3\). When
\(|Z_i|=1\), its order condition also holds:
\(d_{Q_i}(u_{i,1})=d_M(u_{i,1})\ge3\) by (7.2), so
\(|V(A_i)|=|V(Q_i)|\ge4\). Hence each component supplies three
admissible \(u_{i,1},u_{i,2}\)-paths; a light family avoids \(c\),
whereas a heavy family may use it.

Suppose first that at most one component is heavy. Apply the two-link
lemma in \(T\) to \[
 \{y_{1,1},y_{1,2}\}
 \quad\text{and}\quad
 \{y_{2,1},y_{2,2}\}.
\] After swapping the labels \(1,2\) for \(Q_2\), and reversing its
rooted path family if necessary, let the two vertex-disjoint links
\(L_j\) join \(y_{1,j}\) to \(y_{2,j}\) for \(j=1,2\). Orient a selected
path \(P_i\) from \(u_{i,1}\) to \(u_{i,2}\). The links \(L_1,L_2\) have
disjoint vertex sets and have no internal vertex among the four
interface vertices. A selected path \(P_i\) lies in \(Q_i\cup\{c\}\).
Since \(Q_1,Q_2\) are distinct components and at most one of the two
path families can use \(c\), the paths \(P_1,P_2\) are vertex-disjoint.
They are also disjoint from both links, which lie in \(T\).

It follows that \[
 P_1+u_{1,2}y_{1,2}+L_2+y_{2,2}u_{2,2}
 +\overleftarrow{P_2}+u_{2,1}y_{2,1}
 +\overleftarrow{L_1}+y_{1,1}u_{1,1}
\] is a simple cycle. If the selected component-path lengths are
\(p_1,p_2\), its length is \[
 p_1+p_2+|E(L_1)|+|E(L_2)|+4.
\] The connector contribution is fixed, so the two-variable form of
Lemma 2.2 gives five admissible cycles.

Now suppose that both components are heavy. Choose one attachment edge
\(u_iy_i\) from each component with \(y_1\ne y_2\). Let \(K_i\) be
obtained from \(Q_i\) by adjoining \(c\) and all edges from \(c\) to
\(D\cap V(Q_i)\). Since \(Z_i\) contains at least two nonroot vertices,
\(c\) has at least two distinct neighbors in \(Q_i\), so \(K_i\) is
2-connected. Since \(c\notin V(J)\), the roots \(c,u_i\) are distinct,
and \((K_i-cu_i)+cu_i\) contains \(K_i\) as a spanning subgraph and is
2-connected. Apply GHLM Theorem 3.1 with \(q=3\) to \(K_i-cu_i\). Every
nonroot outside \(D\) has degree at least four, and every nonroot in
\(D\) gains its edge to \(c\). Thus each side supplies three admissible
\(c,u_i\)-paths.

A selected \(c,u_i\)-path lies in \(Q_i\cup\{c\}\). Since \(Q_1,Q_2\)
are distinct components, the two selected paths meet exactly at their
common endpoint \(c\). Choose a \(y_1,y_2\)-path \(L\) in \(T\). It is
disjoint from both exterior paths, and \(y_1\ne y_2\). Adjoining
\(u_1y_1\), \(L\), and \(y_2u_2\) therefore produces a simple cycle.

If the two selected path lengths are \(p_1,p_2\), its length is \[
 p_1+p_2+|E(L)|+2.
\] The two-variable form of Lemma 2.2 again gives five admissible
cycles. Both cases are impossible; hence \(M\) is connected. \(\square\)

Since \(M\) is nonempty and connected, it has itself as its unique
component; Lemma 7.1 therefore shows that \(M\) is 2-connected.

\subsubsection{7.3. The final residue
contradiction}\label{the-final-residue-contradiction}

The selected vertices \(x,y\in T\) satisfy
\(d_T(x),d_T(y)\le3<\delta(J)\). Since \(T\) is induced, each has a
neighbor outside \(T\). Choose \(x',y'\in M\), respectively adjacent to
\(x,y\); then \(x'\ne y'\) by (7.1).

\Needspace{6\baselineskip}

Define the exterior exceptional set and root-deleted graph by \[
 Z=D\cap\bigl(V(M)\setminus\{x',y'\}\bigr),
 \qquad
 A=M-x'y'.
\] By construction \(x'y'\notin E(A)\). The graph \(A+x'y'\) is
2-connected because it contains \(M\) as a spanning subgraph. Deleting
the root edge affects only the root degrees, so (7.2) gives the
hypotheses of Lemma 3.3 with \(q=3\). Moreover, \(d_M(x')\ge3\) by
(7.2), so \(|V(A)|=|V(M)|\ge4\). We obtain three admissible
\(x',y'\)-paths whose vertices lie in \(V(M)\cup\{c\}\). Write their
lengths as \[
 a,\ a+d,\ a+2d,
 \qquad d\in\{1,2\}.
\]

\Needspace{12\baselineskip}

Let \(\mathcal R\subseteq\mathbb Z/5\mathbb Z\) be the three distinct
residues of the selected \(x,y\)-paths in \(T\). For a lifted
\(x',y'\)-path \(P\), adjoining \(xx'\) and \(y'y\) produces an
\(x,y\)-path whose internal vertices lie in \(V(M)\cup\{c\}\). Since \[
 V(T)\cap\bigl(V(M)\cup\{c\}\bigr)=\varnothing,
\] this exterior path meets every selected core \(x,y\)-path exactly at
the distinct vertices \(x,y\). Their union is therefore a simple cycle.
If \(P\) is the \(i\)-th lifted path and the core path has residue
\(r\in\mathcal R\), its length is congruent to \[
 a+id+2+r\pmod5,
 \qquad 0\le i\le2.
\] Lemma 2.3, applied with \(a\) replaced by \(a+2\), shows that one of
these nine cycles has length divisible by five. This cycle lies in
\(B\), contradicting the standing assumption that \(B\) contains no such
cycle. Hence the counterexample \(G_0\) cannot exist.

\subsection{8. Completion}\label{completion}

The preceding contradiction proves the main result.

\textbf{Theorem.} Every nonempty finite simple graph of minimum degree
at least five contains a simple cycle whose length is divisible by five.

Combining this theorem with the cases \(k=3\) and \(k=4\), proved by
Chen and Saito\textsuperscript{\hyperref[ref-chen-saito-1994]{{[}2{]}}}
and Dean, Lesniak, and Saito
\textsuperscript{\hyperref[ref-dean-lesniak-saito-1993]{{[}6{]}}},
respectively, and the result of Luo, Ma, and Zhao for every
\(k\ge6\)\textsuperscript{\hyperref[ref-luo-ma-zhao-2026]{{[}8{]}}}
gives the full conjecture:

\textbf{Corollary (Dean's conjecture).} For every integer \(k\ge3\),
every nonempty finite simple graph of minimum degree at least \(k\)
contains a simple cycle whose length is divisible by \(k\).

\newpage

\subsection{References}\label{references}

\begingroup
\raggedright

\begin{enumerate}
\def\labelenumi{\arabic{enumi}.}
\item
  \phantomsection\label{ref-bglp-2026}{}Y. Bai, A. Grzesik, B. Li, and
  M. Prorok, ``Cycle lengths in graphs of given minimum degree,''
  \emph{Journal of Combinatorial Theory, Series B} \textbf{180} (2026),
  111--150.
  \href{https://doi.org/10.1016/j.jctb.2026.06.003}{doi:10.1016/j.jctb.2026.06.003};
  \href{https://arxiv.org/abs/2511.03085v2}{arXiv:2511.03085v2}.
\item
  \phantomsection\label{ref-chen-saito-1994}{}G. T. Chen and A. Saito,
  ``Graphs with a cycle of length divisible by three,'' \emph{Journal of
  Combinatorial Theory, Series B} \textbf{60} (1994), no. 2, 277--292.
  \href{https://doi.org/10.1006/jctb.1994.1019}{doi:10.1006/jctb.1994.1019}.
\item
  \phantomsection\label{ref-chiba-ota-yamashita-2023}{}S. Chiba, K. Ota,
  and T. Yamashita, ``Minimum degree conditions for the existence of a
  sequence of cycles whose lengths differ by one or two,'' \emph{Journal
  of Graph Theory} \textbf{103} (2023), no. 2, 340--358.
  \href{https://doi.org/10.1002/jgt.22921}{doi:10.1002/jgt.22921};
  \href{https://arxiv.org/abs/2008.09783v1}{arXiv:2008.09783v1}.
\item
  \phantomsection\label{ref-choi-chu-kim-park-2026}{}I. Choi, H. Chu, R.
  Kim, and B. Park, ``Existence of cycles of length divisible by 3 or
  4,'' arXiv preprint (2026).
  \href{https://doi.org/10.48550/arXiv.2605.02731}{doi:10.48550/arXiv.2605.02731};
  \href{https://arxiv.org/abs/2605.02731v1}{arXiv:2605.02731v1}.
\item
  \phantomsection\label{ref-dean-1991}{}N. Dean, ``Which graphs are
  pancyclic modulo \(k\)?'' in \emph{Graph Theory, Combinatorics, and
  Applications}, vol.~1, edited by Y. Alavi, G. Chartrand,
  \mbox{O. R. Oellermann}, and A. J. Schwenk, Wiley-Interscience, New
  York, 1991, pp.~315--326.
  \href{https://archive.org/details/graphtheorycombi0000inte}{\mbox{ISBN 978-0-471-53245-3}}.
\item
  \phantomsection\label{ref-dean-lesniak-saito-1993}{}N. Dean, L.
  Lesniak, and A. Saito, ``Cycles of length 0 modulo 4 in graphs,''
  \emph{Discrete Mathematics} \textbf{121} (1993), nos. 1--3, 37--49.
  \href{https://doi.org/10.1016/0012-365X(93)90535-2}{doi:10.1016/0012-365X(93)90535-2}.
\item
  \phantomsection\label{ref-gao-huo-liu-ma-2022}{}J. Gao, Q. Huo, C.-H.
  Liu, and J. Ma, ``A unified proof of conjectures on cycle lengths in
  graphs,'' \emph{International Mathematics Research Notices}
  \textbf{2022} (2022), no. 10, 7615--7653.
  \href{https://doi.org/10.1093/imrn/rnaa324}{doi:10.1093/imrn/rnaa324};
  \href{https://arxiv.org/abs/1904.08126v3}{arXiv:1904.08126v3}.
\item
  \phantomsection\label{ref-luo-ma-zhao-2026}{}Y. Luo, J. Ma, and Z.
  Zhao, ``Dean's conjecture and cycles modulo \(k\),'' arXiv preprint
  (2026).
  \href{https://doi.org/10.48550/arXiv.2601.13552}{doi:10.48550/arXiv.2601.13552};
  \href{https://arxiv.org/abs/2601.13552v1}{arXiv:2601.13552v1}.
\end{enumerate}

\endgroup

\end{document}
